\documentclass[10pt,a4paper]{amsart}
\usepackage[margin=19mm]{geometry}
\usepackage{amsmath,amssymb,mathtools}
\usepackage[scr=rsfs,cal=euler]{mathalfa}
\usepackage{xfrac}
\usepackage[unicode,hidelinks,bookmarks=false]{hyperref}
\newcommand{\ie}{i.e.\ }
\newcommand{\cf}{cf.\ }
\newcommand{\resp}{respectively\ }
\newcommand{\Subsection}{\subsection}
\providecommand{\nobreakdash}{\nobreak\hbox{-}}
\setlength{\emergencystretch}{2em}
\multlinegap0pt
\usepackage{bm}
\usepackage{bbm}
\usepackage{dsfont}
\usepackage{xspace}
\usepackage{cancel}
\usepackage{mathtools}
\usepackage{amsthm}
\makeatletter
\@ifundefined{Remark}{\newtheorem{Remark}{Remark}}{}
\makeatother
\title[Relations between values of arithmetic Gevrey series, and ap]{Relations between values of arithmetic Gevrey series, and applications to values of the Gamma function}
\author{St\'ephane Fischler, Tanguy Rivoal}
\date{Source: arXiv:2301.13518v1 (CC BY 4.0)}
\begin{document}
\maketitle

\noindent\emph{Source and licence.} Relations between values of arithmetic Gevrey series, and applications to values of the Gamma function. Version arXiv:2301.13518v1 (CC BY 4.0), distributed under the Creative Commons Attribution 4.0 International licence, as recorded in the arXiv licence metadata for this exact version and shown on the abstract page https://arxiv.org/abs/2301.13518v1. Full rights notice: licenses/arxiv-2301.13518-CC-BY-4.0.txt.\par\smallskip

\noindent\textit{Editorial scope.} Counted unit: the Remark whose source label is \texttt{None} in the article (source0.tex lines 513--514, source label \texttt{None}), delivered together with the article's own complete original proof of it (source0.tex lines 516--572, one \texttt{proof} environment of 57 lines). No mathematics is written, repaired or re-derived: statement and proof are the article's own text line for line. Required context retained uncounted: none. Every definition, display and label of those lines is retained unchanged. Omitted from this package and not redistributed: 1--512, 515--515, 573--772. Nothing inside a retained excerpt is omitted or altered: every retained body is delivered exactly as the article's lines are written, so no omission receipt of any kind accompanies this package. Citations: the retained excerpts carry no citation command at all, so no bibliography entry of the article is transcribed and no reference is redistributed. Reference adaptation: none was needed and none was made; every reference inside the delivered unit resolves inside it, so no unresolved reference or citation is produced. Printed slips kept literal: no printed word, symbol, hypothesis or number of the counted statement or of its proof was altered. Layout and class: the article's own class is \texttt{article} with packages including a4wide, amsmath, amssymb, amsthm, inputenc, graphicx, babel, xcolor, ulem; the wrapper is the lane's pinned \texttt{amsart} editorial wrapper at 10pt on A4, which restates the theorem-like environments the retained text needs and adds the article's own macro definitions that the retained text needs (none), with extra wrapper packages (bm, bbm, dsfont, xspace, cancel, mathtools). The delivered numbering is the unit's own. Assets: no retained span contains a figure environment, a graphics-inclusion command, a PDF-page inclusion, a table or a TikZ picture; no image file is copied into this package, the package declares no asset dependency and no image file is redistributed with it. Licence: version arXiv:2301.13518v1 of this article is distributed under the Creative Commons Attribution 4.0 International licence, as recorded in the arXiv licence metadata for this exact version and shown on the abstract page https://arxiv.org/abs/2301.13518v1. A full rights notice is delivered at licenses/arxiv-2301.13518-CC-BY-4.0.txt. Characters: every retained excerpt is pure ASCII, so the delivered engine, XeLaTeX with its default Latin Modern text font, reports no missing character.
\begin{Remark} Part $(i)$ of the lemma is false if $a\in \mathbb N^{*}$ because $1, e^z, e^zE_{a,1}(-z)$ are $\mathbb Q(z)$-linearly dependent in this case (see the proof of Part $(ii)$ below).
\end{Remark}

\begin{proof} $(i)$ For simplicity, we set  $\psi_s(z):=e^zE_{a,s}(-z)$. We proceed by induction on $s\ge 1$. Let us first prove the case $s=1$. The derivative of $\psi_1(z)$ is $(1+(z-a)\psi_1(z))/z$. Let us assume the existence of a relation $\psi_1(z)=u(z)e^z+v(z)$ with $u,v\in \mathbb C(z)$ (a putative relation $U(z)+V(z)e^z+W(z)\psi_1(z)=0$ forces $W\neq 0$ because $e^z\notin \mathbb C(z)$). Then after differentiation of both sides, we end up with 
$$
\frac{1+(z-a)\psi_1(z)}{z}=\big(u(z)+u'(z)\big)e^z+v'(z).
$$
Hence, 
$$
\frac{1+(z-a)\big(u(z)e^z+v(z)\big)}{z}=\big(u(z)+u'(z)\big)e^z+v'(z).
$$
Since $e^z\notin \mathbb C(z)$, the function $v(z)$ is a rational solution of the differential equation $zv'(z)=(z-a)v(z)+1$: $v(z)$ cannot be identically 0, and it cannot be a polynomial (the degrees do not match on both sides). It must then have a pole at some point $\omega$, of order $d\ge 1$ say. We must have $\omega=0$ because otherwise the order of the pole at $\omega$ of $zv'(z)$ is $d+1$ while the order of the pole of $(z-a)v(z)+1$ is at most $d$. Writing $v(z)=\sum_{n\ge -d} v_nz^n$ with $v_{-d}\neq 0$ and comparing the term in $z^{-d}$ of $zv'(z)$ and $(z-a)v(z)+1$, we obtain that $d=a$. This forces $a$ to be an integer $\ge 1$, which is excluded. Hence, $1, e^z, e^zE_{a,1}(-z)$ are $\mathbb C(z)$-linearly independent.
\medskip

Let us now assume that the case $s-1\ge 1$ holds. Let us assume the existence of a relation over $\mathbb C(z)$
\begin{equation} \label{eq:psisz}
    \psi_s(z)=v(z)+u_0(z)e^z+\sum_{j=1}^{s-1} u_j(z)\psi_j(z).
\end{equation}
(A putative relation $V(z)+U_0(z)e^z+\sum_{j=1}^{s}U_j(z)\psi_j(z)=0$ forces $U_s\neq 0$ by the induction hypothesis). Differentiating \eqref{eq:psisz} and because 
 $\psi_j'(z)=(1-\frac{a}{z})\psi_j(z)+\frac{1}{z}\psi_{j-1}(z)$
for all $j\ge 1$ (where we have let $\psi_0(z)=1 $), we have 
\begin{multline} \label{eq:psisz2}
  A(z)\psi_s(z)+\frac{1}{z}\psi_{s-1}(z) =
   v'(z)+\big(u_0(z)+u'_0(z)\big)e^z
   +\sum_{j=1}^{s-1}u'_j(z)\psi_j(z)\\
+\sum_{j=1}^{s-1}u_j(z)\big(A(z)\psi_j(z)+\frac{1}{z}\psi_{j-1}(z)\big), 
\end{multline}
where $A(z):=1-a/z$. 
Substituting the right-hand side of \eqref{eq:psisz} for $\psi_s(z)$ on the left-hand side of \eqref{eq:psisz2}, we then deduce that
\begin{multline*}
    v'(z)-A(z)v(z)+\big(u'_0(z)
    +(1-A(z)) u_0(z) \big)e^z\\
    +\frac{1}{z}(z-a)u_1(z)\psi_1(z)+\sum_{j=1}^{s-1}u'_j(z)\psi_{j}(z)+\frac{1}{z}\sum_{j=1}^{s-1} u_j(z)\psi_{j-1}(z)
-\frac{1}{z}\psi_{s-1}(z) =0.
\end{multline*}
This is  a  non-trivial $\mathbb C(z)$-linear relation between 
$ 1, e^z, \psi_1(z), \psi_2(z), \ldots, \psi_{s-1}(z) $ because the coefficient of $\psi_{s-1}(z)$ is $u_{s-1}'(z)-1/z$ and it is  not identically 0 because $u_{s-1}'(z)$ cannot have a pole of order $1$. But by the induction hypothesis, we know that such a relation is impossible. 

\medskip

$(ii)$ The proof can be done by induction on $s\ge 2$ similarily. In the case $s=2$, assume the existence of a relation $\psi_2(z)=u(z)e^z+v(z)$ with $u(z), v(z)\in \mathbb C(z)$. By differentiation, we obtain
$$
\Big(1-\frac{a}{z}\Big)\psi_2(z)=-\frac{1}{z}\psi_1(z)+\big(u(z)+u'(z)\big)e^z+v'(z).
$$
By induction on $a\ge 1$, we have $\psi_1(z)=(a-1)!e^z/z^{a}+w(z)$ for some $w(z)\in \mathbb Q(z)$. Hence, we have 
$$
\Big(1-\frac{a}{z}\Big)u(z)=-\Big(\frac{(a-1)!}{z^{a+1}}+1\Big)u(z)+u'(z)
$$
which is not possible. Let us now assume that the case $s-1\ge 2$ holds, as well as the existence of a relation over $\mathbb C(z)$
\begin{equation} \label{eq:psisznew}
    \psi_s(z)=v(z)+u_0(z)e^z+\sum_{j=2}^{s-1} u_j(z)\psi_j(z).
\end{equation}
We proceed exactly as above by differentiation of both sides of \eqref{eq:psisznew}. Using the relation  $\psi_j'(z)=(1-\frac{a}{z})\psi_j(z)+\frac{1}{z}\psi_{j-1}(z)$
for all $j\ge 2$ and the fact that $\psi_1(z)=(a-1)!e^z/z^{a}+w(z)$, we obtain a relation 
$
\widetilde{v}(z)+\widetilde{u}_0(z)e^z+\sum_{j=2}^{s-1} \widetilde{u}_j(z)\psi_j(z) =0 
$
where $\widetilde{u}_{s-1}(z)=u_{s-1}'(z)-1/z$ 
cannot be identically 0. The induction hypothesis rules out the existence of such a relation.
\end{proof}

\end{document}
